\chapter{Literature Review}

\section*{2.0 Overview}
\addcontentsline{toc}{section}{2.0 Overview}
This chapter presents a critical review of relevant literature in human intention recognition, optical gesture analysis, and autonomous robot navigation in collaborative workspaces. Existing research establishes that autonomous mobile systems can recognize gestures, infer human motion trajectories, and adapt navigation behaviors to enhance interaction safety. However, a significant proportion of these systems operate under controlled laboratory conditions, depend on specialized high-power sensing hardware, or decouple intention recognition from mobile navigation. Consequently, this study builds directly upon existing theoretical and empirical frameworks by integrating vision-based human motion tracking and geometric hand gesture classification into an edge-deployed, mobile robotic platform operating in a mapped indoor environment.

\section{Theoretical Review}
This research is theoretically grounded in the principles of Human-Robot Interaction (HRI), human intention theory, geometric landmark kinematics, and spatial mobile robotics. Human intention theory posits that human goal-directed actions contain early predictive cues embedded within movement trajectories, forming the computational foundation for anticipatory robotic response systems. Within contemporary Human-Robot Interaction paradigms, collaborative mobile robots are required to move beyond purely reactive obstacle avoidance toward proactive, context-aware interaction within human-shared operational domains.

The theoretical formulation of gesture and motion recognition in this study is based on the premise that human hand anatomical configurations can be modeled geometrically and topologically from single-frame landmark distributions. Static gesture classification relies on the principle that instantaneous spatial arrangements of hand joints provide sufficient discriminative features to distinguish discrete operational commands within a fixed vocabulary. Conversely, human movement trajectory prediction inherently requires temporal context, as instantaneous positions cannot reveal directional intent or velocity. This distinction necessitates a sequential modeling approach—specifically Long Short-Term Memory (LSTM) recurrent neural networks—to capture temporal dependencies across consecutive pose observations over a sliding temporal window.

Spatial mapping and localization theory provides the mathematical framework for continuous autonomous navigation between human interaction events. Simultaneous Localization and Mapping (SLAM) constructs an occupancy-grid representation of the physical environment, while probabilistic Monte Carlo algorithms estimate the robot's coordinate pose relative to the global reference frame. Finally, multi-modal perceptual fusion theory demonstrates that combining explicit command channels (hand gestures) with implicit contextual cues (kinematic motion trajectories) mitigates command ambiguity and enhances operational robustness. Together, these theoretical frameworks provide the formal basis for developing a unified robotic system capable of recognizing, predicting, and responding to human intentions on physical hardware.

\section{Review of Related Works}

\subsection{Review One: Real-Time Intention Prediction via Upper-Limb Kinematics}

\subsubsection{Research Topic and Authors}
\textit{Real-time Feasibility of a Human Intention Method Evaluated Through a Competitive Human--Robot Reaching Game} (Tsitos et al., 2022) \cite{tsitos2022}.

\subsubsection{Objectives}
The primary objective of Tsitos et al. was to develop and evaluate an edge-feasible human intention prediction framework using a single RGB-D camera and the OpenPose skeletal tracking library. The authors aimed to integrate this predictive module into a standard ROS-based robotic control pipeline to govern arm trajectories in real time. Furthermore, the study investigated whether an industrial manipulator could react sufficiently early during competitive reaching tasks based on anticipatory kinematics rather than post-motion feedback.

\subsubsection{Methodology}
The authors implemented an experimental pipeline combining real-time human pose estimation with machine learning classification. An RGB-D camera captured upper-body human motion, from which OpenPose extracted 2D and 3D wrist joint coordinates during reaching movements. Raw spatial coordinates were digitally filtered to eliminate tracking jitter and segmented into distinct kinematic movement phases. Feature vectors were subsequently constructed from velocity and position profiles to train multiple classification models, including Support Vector Machines (SVM), Logistic Regression, Decision Trees, and Naive Bayes classifiers. The finalized predictive pipeline was interfaced directly into a ROS environment to control a 6-DOF Universal Robots UR3 manipulator during competitive reaching tasks.

\subsubsection{Strengths}
The investigation demonstrates that accurate intention prediction can be achieved using a streamlined, cost-effective setup comprising a single monocular depth sensor. By restricting the feature space to minimal wrist joint coordinates, the system sustained low-latency inference suitable for time-critical robotic interaction. The physical deployment on an industrial UR3 manipulator provided empirical validation on real hardware rather than simulated kinematics. Additionally, the experimental identification of a strict human reaction time threshold emphasized the critical role of processing latency in collaborative robotics.

\subsubsection{Limitations}
A primary limitation of the study is that the evaluation was confined to a static tabletop reaching game without mobile base movement. The human operator and the robotic arm did not share a contiguous mobile floor space, preventing the evaluation of spatial navigation or dynamic obstacle avoidance. Furthermore, the system relied exclusively on implicit reaching motion, offering no explicit gesture command channel through which the human could issue deliberate supervisory directives.

\subsubsection{Research Gap}
While Tsitos et al. demonstrated that kinematic intention prediction can operate within strict latency bounds, their architecture is restricted to a stationary manipulator arm. The robot and human never navigate a shared physical floor environment, and the human possesses no mechanism to command the robot through deliberate, explicit gestures. This leaves a distinct engineering gap between fast, stationary intention inference and its deployment onto an autonomous mobile platform capable of accepting both explicit gesture commands and navigating physical spaces.

\subsection{Review Two: 3D Skeletal Tracking and Adaptive Gesture Classification}

\subsubsection{Research Topic and Authors}
\textit{3D Gesture Recognition and Adaptation for Human--Robot Interaction} (Mahmud et al., 2022) \cite{mahmud2022}.

\subsubsection{Objectives}
The core objective of Mahmud et al. was to develop a three-dimensional hand gesture recognition and adaptation architecture capable of identifying both pointing and dynamic gestures in real time using depth sensing. The study sought to apply these classified gestures directly to robot navigation within human-robot interaction domains. Additionally, the authors aimed to establish a semi-supervised gesture adaptation mechanism that allowed the robotic system to incorporate novel user-defined gestures dynamically.

\subsubsection{Methodology}
The system employed a Microsoft Kinect depth sensor to capture 3D skeletal joint coordinates across the operator's upper torso and hands. Pointing vectors were computed geometrically using 3D line equations projected onto the horizontal X--Z plane, while dynamic gestures were modeled from sequential joint trajectories over time. Spatial joint coordinates were normalized relative to a torso reference frame to maintain geometric consistency across different user physical builds. These feature representations were evaluated across three distinct machine learning classifiers: Hidden Markov Models (HMM), Support Vector Machines (SVM), and Convolutional Neural Networks (CNN). The navigation pipeline was evaluated within a simulated indoor environment using a dataset of 3,600 gesture instances collected from 24 participants across multiple age demographics.

\subsubsection{Strengths}
The utilization of three-dimensional skeletal coordinate data substantially enhanced tracking robustness against ambient lighting shifts and visual background clutter compared to conventional 2D vision techniques. The simultaneous integration of pointing vectors and dynamic gestures provided an expressive, natural interaction vocabulary for human operators. Moreover, the semi-supervised adaptation mechanism granted flexibility by allowing the system to update gesture vocabularies without requiring complete model retraining from scratch. The extensive demographic evaluation across multiple age brackets demonstrated practical insights into user adaptability across varying motor proficiencies.

\subsubsection{Limitations}
Despite these operational strengths, the architecture exhibits severe hardware and environmental constraints. The entire tracking pipeline depends strictly on the proprietary Microsoft Kinect hardware, restricting deployment to devices supporting heavy depth-sensor drivers. Furthermore, the navigation component was validated exclusively within a simulated software environment, leaving real-world sensor noise, wheel slip, and embedded latency unaddressed. The pointing estimation relied exclusively on right-hand joint configurations, which restricts natural interaction for left-handed operators or bi-manual interaction.

\subsubsection{Research Gap}
Mahmud et al. established a robust multi-modal gesture classification pipeline, but their navigation control was evaluated only in simulation and relied fundamentally on a high-power Kinect depth sensor. In contrast, non-visual approaches such as surface electromyography (sEMG) explored by Zafar et al. \cite{zafar2023} eliminate visual occlusion issues, but impose the significant operational burden of wearable electrodes on the operator. This creates a compelling need for a completely device-free, monocular camera-based framework that performs gesture classification and spatial navigation entirely on edge hardware in real physical environments.

\subsection{Review Three: Intention-Aware Navigation and Costmap Integration}

\subsubsection{Research Topic and Authors}
\textit{Safe and Efficient Motion Planning for Material Transportation Robots Considering Intention Prediction of Obstacles} (Li et al., 2023) \cite{li2023}.

\subsubsection{Objectives}
The primary objective of Li et al. was to formulate an intention-aware motion planning framework for autonomous transportation robots operating in congested industrial construction sites. The authors aimed to extend conventional static obstacle avoidance by predicting whether detected obstacles (specifically human workers) were likely to yield or move away from the robot's projected trajectory. Furthermore, the study sought to integrate these intention predictions directly into ROS2 Navigation2 (Nav2) costmaps without modifying the underlying path planning algorithms.

\subsubsection{Methodology}
The authors developed a multi-sensor perception architecture combining 2D LiDAR scanning with camera-based convolutional object detection to classify site obstacles into distinct categories, including human workers, transport carts, machinery, and static materials. A predictive intention model based on Rectified Linear Unit (ReLU) activation functions estimated the probability of obstacle clearance based on semantic class, relative velocity, and Euclidean distance. When an obstacle was predicted to clear the path, the global costmap was dynamically updated to prevent unnecessary path replanning, while the local costmap maintained active collision safety. The complete navigation pipeline was implemented within the ROS2 Nav2 framework and evaluated using Isaac Sim on an autonomous Carter mobile platform across simulated construction environments.

\subsubsection{Strengths}
The central engineering strength of Li et al.'s framework is the dynamic modification of standard global costmaps based on predicted obstacle behavior, which significantly reduces unnecessary robot stops and inefficient path deviations. The multi-modal combination of LiDAR distance arrays and visual semantic detection provided rich spatial understanding in complex layouts. Furthermore, because the architecture interfaces directly with standard Nav2 costmap layers, it remains readily portable across different mobile base platforms without altering the core navigation solvers.

\subsubsection{Limitations}
A significant limitation is that the entire evaluation was conducted inside the NVIDIA Isaac Sim virtual simulation environment without physical robot deployment. As a consequence, physical motor actuation delays, real sensor dropouts, and embedded memory constraints were not verified on constrained edge processors. Furthermore, the intention recognition was entirely passive; the system possessed no explicit communication channel allowing human workers to deliberately direct, pause, or interact with the robot through gestures.

\subsubsection{Research Gap}
Li et al. demonstrated the architectural utility of modifying Nav2 costmaps using inferred human intentions, but their framework remains strictly passive and unverified on physical hardware. The robot infers whether a worker might clear the path, but the worker cannot issue an explicit command (such as an emergency stop or a directional waypoint) through natural visual gestures. This reveals a distinct research gap: combining an active, explicit gesture command channel with autonomous navigation on a physically deployed, edge-only mobile platform.

\subsection{Synthesis of Reviewed Literature}
A comprehensive analysis of the literature reveals a clear technological divide across current research. Systems that achieve exceptional intention prediction accuracy (such as Tsitos et al. \cite{tsitos2022}) operate predominantly on stationary robotic arms decoupled from mobile navigation. Conversely, frameworks that implement intention-aware mobile navigation (such as Li et al. \cite{li2023} and Yu et al. \cite{yu2024}) rely solely on passive obstacle classification in virtual simulation environments, lacking explicit human-command interfaces. Furthermore, recent lightweight gesture architectures (such as Muhtadin et al. \cite{muhtadin2025}) demonstrate edge-feasible MediaPipe classification on ROS2, but validate their models strictly on stationary UR5 industrial arms without mobile chassis or SLAM-based spatial mapping.

This study directly resolves this collective research gap by engineering a unified, physical mobile robotic platform that integrates active, scale-invariant hand gesture recognition, temporal motion prediction, and LiDAR SLAM mapping natively onto an embedded edge computing architecture.

\begin{landscape}
\section{Comparative Summary of Additional Reviewed Literature}

\footnotesize
\begin{longtable}{|
>{\RaggedRight\arraybackslash}p{3.2cm}|
>{\RaggedRight\arraybackslash}p{3.5cm}|
>{\RaggedRight\arraybackslash}p{5.2cm}|
>{\RaggedRight\arraybackslash}p{3.8cm}|
>{\RaggedRight\arraybackslash}p{3.8cm}|
>{\RaggedRight\arraybackslash}p{3.8cm}|
}
\caption{Comparative Summary of Reviewed Literature in Gesture, Intention, and Robot Navigation}
\label{tab:lit_review_landscape}\\
\hline
\textbf{Authors / Year} & \textbf{Research Topic} & \textbf{Methodology} & \textbf{Limitations} & \textbf{Research Gap} & \textbf{Proposed Improvement} \\
\hline
\endfirsthead

\hline
\textbf{Authors / Year} & \textbf{Research Topic} & \textbf{Methodology} & \textbf{Limitations} & \textbf{Research Gap} & \textbf{Proposed Improvement} \\
\hline
\endhead

Xinbo Yu et al. (2024) \cite{yu2024} &
Real-Time Trajectory Planning and Obstacle Avoidance for Human--Robot Co-Transporting &
Employs a physical slider-sensor mechanism and 2D LiDAR with a multiscale perception model to estimate human pulling motion and plan safe co-transporting trajectories. &
Relies on physical contact via mechanical sliders; indirect intention inference easily perturbed by mechanical friction in dynamic environments. &
Determines human intention through mechanical sliders rather than direct, non-contact optical tracking of human motion and gestures. &
Eliminate mechanical contact interfaces by deploying touchless monocular vision tracking combined with planar LiDAR spatial mapping. \\
\hline

Javier Laplaza et al. (2022) \cite{laplaza2022} &
Context and Intention for 3D Human Motion Prediction in Handover Tasks &
Attention-based recurrent deep learning model that predicts 3D human reaching motion using motion history, robot spatial pose, and environmental object context. &
Experiments restricted to stationary tabletop handover benchmarks and synthetic datasets; does not generalize to mobile navigation spaces. &
Confined to localized object handover tasks; lacks integration with mobile platform locomotion or autonomous navigation middleware. &
Extend predictive modeling to mobile spatial navigation, coupling short-horizon motion prediction with dynamic vehicle control. \\
\hline

Yongshi Liang and Pai Zheng (2024) \cite{liang2024} &
Human Intention Prediction with Visual-Language Model &
Utilizes multi-modal Visual-Language Models (VLMs) to infer human intention from combined visual video streams and linguistic prompts in shared spaces. &
Extremely high computational footprint requiring multi-GPU server infrastructure; tested in pedestrian traffic without robotic control integration. &
Focuses purely on high-level semantic intention prediction; lacks real-time edge execution or physical robotic middleware integration. &
Replace computationally heavy VLMs with lightweight, engineered geometric feature extractors deployable on embedded ARM processors. \\
\hline

Muhtadin et al. (2025) \cite{muhtadin2025} &
Hand Gesture Recognition for Collaborative Robots Using Lightweight Deep Learning in Real-Time Robotic Systems &
MediaPipe hand landmark extraction feeding a lightweight artificial neural network, integrated into ROS2 to control a UR5 industrial manipulator arm in real time. &
Deployed strictly on a stationary 6-DOF robotic arm; no mobile base, no SLAM mapping, and no spatial obstacle perception. &
Demonstrates an edge-feasible gesture pipeline, but does not extend to mobile autonomous robots or multi-modal spatial navigation. &
Bridge the lightweight gesture recognition pipeline with mobile chassis kinematics, 2D LiDAR SLAM, and centralized decision state machines. \\
\hline

\end{longtable}
\normalsize
\end{landscape}
